\chapter{Ein uraltes Kirchengebet}
\quellseite{8}

Um das Jahr 96 nach Christo richtete im Auftrag der Gemeinde zu Rom einer ihrer Prediger, Namens Clemens, ein Schreiben an die Christengemeinde zu Corinth, um eine dort ausgebrochene Zwistigkeit durch brüderlichen Zuspruch beizulegen. Dieses Sendschreiben wurde um seines köstlichen Inhalts willen in den letzten 250 Jahren mehr als vierzigmal griechisch und auch ein dutzendmal in deutscher Uebersetzung gedruckt. Immer aber fehlten einige Blätter am Schluß, und kein Gelehrter konnte sagen, was auf denselben gestanden habe, kein Gelehrter wagte auch zu hoffen, daß man, nachdem alle Bibliotheken Europas durchsucht worden waren, das Fehlende noch finden möchte. Und doch, was niemand gedacht hatte, geschah. Vor etwa fünfundzwanzig Jahren fand man in einer Büchersammlung zu Constantinopel eine vollständige Handschrift dieses merkwürdigen Sendschreibens. Da stellte es sich heraus, daß der Brief noch mehrere Capitel länger gewesen war, als man bis dahin gewußt hatte, und, was am werthvollsten ist, daß er mit einer Fürbitte und einem Gebet schließt, das uns eine Vorstellung davon gibt, wie in jener ältesten Zeit, dreißig Jahre nach dem Tod eines Paulus und Petrus, das sonntägliche Kirchengebet in der Christengemeinde zu Rom lautete. Besonders beachtenswerth ist dabei, daß die Fürbitte für die Obrigkeit, für den heidnischen Kaiser, vor allem hervortritt. Und dabei muß man bedenken, daß der römische Kaiser, für den in dieser Weise gebetet wurde, kein anderer als Domitian war, ein bitterer Christenfeind und Verfolger, „ein Stück Nero an Grausamkeit“, wie ihn ein späterer christlicher Kirchenvater nannte. Dieses uralte christliche Gemeindegebet lautet also:

„Herr, laß erscheinen über uns dein Antlitz zum Guten im Frieden, damit wir bedeckt werden durch deine starke Hand und errettet werden von allen Sünden durch deinen hohen Arm. Und rette uns von denen, die uns ungerecht hassen. Gib Eintracht und Frieden uns und allen, die auf Erden wohnen, wie du sie unsern Vätern gegeben hast, da sie dich anriefen in Glauben und Wahrheit, und hilf uns, die wir unterthan sind deinem allherrschenden und herrlichen Namen. Unsern Obrigkeiten aber und Herrschern auf der Erde hast du, Herr, gegeben die Gewalt ihres Königthums vermöge deiner herrlichen und unausforschlichen Macht, damit wir uns ihnen unterordnen, in nichts widerstrebend deinem Willen.
Gib ihnen Gesundheit, Frieden, Eintracht, Wohlergehen, damit sie die von dir ihnen gegebene Herrschaft ausüben ohne Anstoß. Lenke ihren Rath nach dem, was vor dir gut und wohlgefällig ist, damit sie in Frieden und Sanftmuth die ihnen gegebene Gewalt verwalten und deine Gnade erlangen. Du bist allein im Stande, diese und noch reichlichere Güter uns zu verleihen; vor dir legen wir unser Bekenntniß ab durch den Hohenpriester und Vorsteher unserer Seelen, Jesus Christus, durch welchen dir sei Ruhm und Herrlichkeit jetzt und von Geschlecht zu Geschlecht und von Ewigkeit zu Ewigkeit! Amen.“
